Sources : France Stratégie, comité d'évaluation des réformes de la fiscalité du capital, rapport final (2023), données POTE 2021 ; Assemblée nationale, rapport n° 3056 (2026), d'après la DGFiP, pour les plus-values mises en report en 2023.\par Note : part reçue par les 0,01 \% de foyers fiscaux au revenu le plus élevé, soit environ 4 000 foyers. Le 1 \% le plus aisé déclare 96 \% des dividendes ; 84 foyers ont mis en report 70 \% des plus-values de 2023.
